and shall first prove 6.10 for @@M0@@ and @@M1@@ (which do satisfy the hypotheses of 6.10). In view of the preceding remark on the filtration of @@M2@@, it suffices for this to show that:
@@M3@@
Now we have @@M4@@; this relation can therefore be written:
@@M5@@
Put @@M6@@; then @@M7@@, and we are reduced to showing:
@@M8@@
Now this follows from 6.9.

To deduce 6.10 from this particular case, we consider the @@M9@@-@@M10@@-homomorphism @@M11@@ which sends @@M12@@ to @@M13@@ and @@M14@@ to @@M15@@ for every @@M16@@ and every @@M17@@ (4.7). This homomorphism is compatible with the @@M18@@-filtrations, and sends @@M19@@ to @@M20@@, and @@M21@@ to @@M22@@, whence the result.

\sgaStableTarget{sga6.e05.prop.6-11}\textbf{Proposition 6.11.} Suppose that @@M23@@ is of characteristic zero. Let @@M24@@ be an integer. Then the functor @@M25@@ is an equivalence of categories from the category of graded @@M26@@-algebras @@M27@@ such that @@M28@@, @@M29@@ for @@M30@@, to the category of augmented @@M31@@-@@M32@@-algebras @@M33@@ such that @@M34@@; a quasi-inverse functor is given by @@M35@@.

Let @@M36@@ be a graded @@M37@@-algebra such that @@M38@@ for @@M39@@, endowed with the filtration associated to its grading; then @@M40@@, hence @@M41@@ (6.6.5). Conversely, if @@M42@@ is an augmented @@M43@@-@@M44@@-algebra such that @@M45@@, then @@M46@@, @@M47@@ for @@M48@@.

We shall define isomorphisms of functors:
@@M49@@

